\PaperPage{15}{2.1 Bernstein rows and mass}
\label{bern:rows}

For this chapter let $u=(u_k)_{k\ge0}$ be an arbitrary real sequence.
Normalization and positivity will be imposed explicitly when needed.
Associate to $u$ the polynomial moment functional
\[
 \mathcal M_u\!\left(\sum_{k=0}^d c_kx^k\right)
       =\sum_{k=0}^d c_ku_k,
 \qquad B_{N,j}(x)=\binom Njx^j(1-x)^{N-j}.
\]
This functional is defined by finite algebra; no representing measure,
continuity assertion, or growth restriction on $u$ is implicit. Applying
it to the Bernstein basis gives the weights
\begin{equation}\label{bern:weights}
 w^u_{N,j}=\mathcal M_u(B_{N,j})
 =\binom Nj\sum_{r=0}^{N-j}(-1)^r\binom{N-j}{r}u_{j+r},
 \qquad 0\le j\le N.
\end{equation}
The superscript will be suppressed when the sequence is fixed. Each
row depends on only $u_0,\ldots,u_N$. Thus every particular weight is
finite data, even though positivity of all rows is an infinite family
of conditions.

With the forward-difference convention $\Delta u_j=u_{j+1}-u_j$,
the inner sum is $(-1)^{N-j}\Delta^{N-j}u_j$. Since the binomial
factor is strictly positive at every in-range index, positivity of
all Bernstein rows is equivalent to nonnegativity of every alternating
finite difference. The row formulation retains the binomial weights
because they organize both conservation of mass and the compatibility
between successive degrees. These two identities will control the
normalization of the eventual representing measure.

\begin{proposition}[Conservation of row mass]\label{bern:mass}
For every real sequence and every $N\ge0$,
\[
 \sum_{j=0}^Nw^u_{N,j}=u_0.
\]
\end{proposition}
\begin{proof}
The binomial theorem gives the polynomial identity
$\sum_jB_{N,j}(x)=(x+(1-x))^N=1$. Apply the linear functional
$\mathcal M_u$. Equivalently, expansion of the defining finite sum
shows that the coefficient of $u_k$ is
\[
 \sum_{j=0}^k\binom Nj(-1)^{k-j}\binom{N-j}{k-j}
 =\binom Nk\sum_{j=0}^k\binom kj(-1)^{k-j}
 =\begin{cases}1,&k=0,\\0,&1\le k\le N.\end{cases}
\]
Here the product of binomial coefficients counts the same successive
choices in two orders. All rearrangements are finite, including when
$N=0$.
\end{proof}

The first rows make the sign convention transparent:
\[
 \begin{gathered}
 (w_{0,0})=(u_0),\qquad
 (w_{1,0},w_{1,1})=(u_0-u_1,u_1),\\
 (w_{2,0},w_{2,1},w_{2,2})=(u_0-2u_1+u_2,2u_1-2u_2,u_2).
 \end{gathered}
\]
If a positive measure already represents $u$, these are integrals of
nonnegative polynomials. Our later construction reverses that implication:
the signs of these finite expressions will produce the measure.
Conservation of mass is the normalization mechanism for that reversal.
In particular, when $u_0=1$, every row has signed mass one regardless
of its individual signs. Its absolute mass can nevertheless exceed one;
that excess, rather than the signed mass, detects cancellation.

\PaperPage{16}{2.2 Exact coarsening and increasing variation}
\label{bern:coarsening-page}

The Bernstein bases at successive degrees satisfy a compatibility
relation. It provides a direct comparison of row variation without
assuming that any weights are positive. For $0\le j\le n$, the two
binomial identities
\[
 (n+1-j)\binom{n+1}{j}=(n+1)\binom nj,
 \qquad
 (j+1)\binom{n+1}{j+1}=(n+1)\binom nj
\]
give an identity between polynomials. Factoring out
$\binom njx^j(1-x)^{n-j}$ from its right-hand side leaves
$(n+1)((1-x)+x)$. Consequently,
\begin{equation}\label{bern:poly-coarsen}
 (n+1)B_{n,j}
 =(n+1-j)B_{n+1,j}+(j+1)B_{n+1,j+1}.
\end{equation}
The endpoint cases $j=0,n$ require no separate convention. Applying
$\mathcal M_u$ proves the exact weighted relation
\begin{equation}\label{bern:coarsen}
 (n+1)w_{n,j}
 =(n+1-j)w_{n+1,j}+(j+1)w_{n+1,j+1}.
\end{equation}
The coefficients become nonnegative after division by $n+1$.
They describe how adjacent entries of a finer row contribute to one
coarser entry. They need not sum to one within that pair: the relevant
conservation law is obtained by summing contributions from the whole
finer row.

\begin{proposition}[Monotonicity of variation]\label{bern:variation}
Define $V_N=\sum_{j=0}^N|w_{N,j}|$. Then
\[
 |u_0|\le V_N\le V_{N+1}\qquad(N\ge0),
 \qquad V_0=|u_0|.
\]
\end{proposition}
\begin{proof}
The first inequality is the triangle inequality applied to row mass.
Taking absolute values in~\eqref{bern:coarsen} and summing gives
\[
 (n+1)V_n\le
 \sum_{j=0}^n\bigl((n+1-j)|w_{n+1,j}|
                  +(j+1)|w_{n+1,j+1}|\bigr)
 =(n+1)V_{n+1}.
\]
For an interior finer-row index $k$, the two coefficients are
$n+1-k$ and $k$, whose sum is $n+1$. The indices $0$ and $n+1$
each occur once, again with coefficient $n+1$. Division by this
positive integer proves monotonicity.
\end{proof}

Equation~\eqref{bern:coarsen} also shows that positivity propagates
downward in degree: if every entry of row $n+1$ is nonnegative,
then so is every entry of row $n$. Iteration gives positivity of
every earlier row from positivity of one later row. This is a
finite implication, with the direction fixed by the coarsening
identity. It supplies no positivity information about rows of higher
degree than those already controlled.

This proof identifies the only possible source of strict growth:
the triangle inequality at an adjacent pair. There is no loss from
reindexing or from normalization. The next page records that loss
exactly, rather than estimating it. This exactness is useful because
a zero bound will force equality at every finite step.

\PaperPage{17}{2.3 Triangle defects and the cancellation identity}
\label{bern:cancellation-page}

Fix a row transition $n\to n+1$. For $0\le j\le n$, set
\[
 a_{n,j}=(n+1-j)w_{n+1,j},\qquad
 b_{n,j}=(j+1)w_{n+1,j+1},\qquad
 D_{n,j}=|a_{n,j}|+|b_{n,j}|-|a_{n,j}+b_{n,j}|.
\]
The quantity $D_{n,j}$ is a nonnegative real number. For real scalars,
it vanishes exactly when $a_{n,j}b_{n,j}\ge0$, including either scalar
being zero. Indeed, if the signs agree, absolute value is additive;
if the signs differ, the defect is twice the smaller absolute value.
Thus the terminology ``cancellation'' has a literal algebraic meaning.
Define the total defect of the transition by
\[
 C_n=\sum_{j=0}^nD_{n,j}\ge0.
\]

\begin{proposition}[Exact variation increment]\label{bern:increment}
Every real sequence satisfies
\begin{equation}\label{bern:increment-eq}
 C_n=(n+1)(V_{n+1}-V_n).
\end{equation}
\end{proposition}
\begin{proof}
Equation~\eqref{bern:coarsen} implies
$|a_{n,j}+b_{n,j}|=(n+1)|w_{n,j}|$.
The coefficient count in the preceding proof gives
$\sum_j(|a_{n,j}|+|b_{n,j}|)=(n+1)V_{n+1}$.
Subtract the sum of the former identities from the latter identity.
Every sum is finite, so this is an equality with no convergence issue.
\end{proof}

Two useful consequences retain the full information of the defect:
\[
 C_n=0\ \Longleftrightarrow\ V_{n+1}=V_n,
 \qquad
 V_N=|u_0|+\sum_{n=0}^{N-1}\frac{C_n}{n+1}.
\]
The second identity follows by telescoping, with an empty sum when
$N=0$. It separates the initial absolute mass from the accumulation of
local cancellation. In particular, a bound on the weighted sum of
defects bounds total variation; an assertion that all defects vanish
fixes total variation exactly.

When $u_0=1$, another exact expression is
\[
 V_N-1=2\sum_{j=0}^N\max\{-w_{N,j},0\}.
\]
To prove it, subtract row mass from absolute mass and use
$|t|-t=2\max\{-t,0\}$ term by term. This makes the distinction
between signed mass and negative mass explicit. One zero transition
does not assert positivity of every row. Positivity at all degrees
will follow from zero defects at all transitions together with the
initial positive mass.

For example, the normalized initial moments $u_0=1$, $u_1=u_2=-1$
give rows $(2,-1)$ and $(2,0,-1)$. Both have variation three, so
$C_1=0$ despite their negative entries. The earlier transition has
$C_0=2$. This example isolates why all transitions, including the
transition out of the degree-zero row, occur in the zero-budget
hypothesis.

The source definitions
\texttt{cancellationDefect} and \texttt{rowCancellation} encode these
same weighted triangle defects~\cite{source}.

\PaperPage{18}{2.4 The zero budget and its equivalent forms}
\label{bern:zero-budget-page}

For a normalized real sequence, a cancellation budget is a pair of
natural numbers $W=(N_0,b)$. Its semantic admissibility means
\[
 \operatorname{Adm}_u(N_0,b)\quad\Longleftrightarrow\quad
 \forall n\ge N_0,\ C_n\le bV_n.
\]
Write $H_{00}(u)$ for $\operatorname{Adm}_u(0,0)$. This notation
temporarily makes the sequence explicit. In the actual application,
$H_{00}$ without an argument always refers to the fixed normalized
logarithmic-derivative moments specified in Part~4.

\begin{theorem}[Finite identities behind the zero budget]\label{bern:h00}
If $u_0=1$, then the following four assertions are equivalent:
\begin{equation}\label{bern:h00-equivalences}
 \begin{aligned}
 H_{00}(u)
 &\ \Longleftrightarrow\ \forall n\ge0,\ C_n=0\\
 &\ \Longleftrightarrow\ \forall N\ge0,\ V_N=1\\
 &\ \Longleftrightarrow\ \forall N\ge0\ \forall j\le N,\ w_{N,j}\ge0.
 \end{aligned}
\end{equation}
\end{theorem}
\begin{proof}
Admissibility gives $C_n\le0$ at every index; nonnegativity of defects
turns this inequality into equality. Conversely, zero defects satisfy
the budget. By~\eqref{bern:increment-eq}, zero defects are equivalent
to equal successive variations. Since $V_0=|u_0|=1$, induction gives
$V_N=1$ at every degree. Constant variation conversely makes every
increment vanish.

For a fixed row, conservation of mass yields
\[
 V_N-1=\sum_{j=0}^N\bigl(|w_{N,j}|-w_{N,j}\bigr).
\]
Each summand is nonnegative. A finite sum of such terms is zero
exactly when every term is zero, and $|t|-t=0$ exactly when $t\ge0$.
This proves both remaining directions, row by row.
\end{proof}

The role of normalization is precise. It supplies a positive initial
mass and identifies minimal variation with one. If $u_0=c>0$, the
same argument replaces one by $c$. For the actual sequence,
normalization $u_0=1$ is already a retained theorem, so no additional
normalization hypothesis is attached to the actual implication.

The pair $(0,0)$ has a finite description, but its admissibility is
the universally quantified assertion displayed above. Checking any
fixed collection of rows proves only the corresponding finite
collection of inequalities. The theorem therefore gives equivalent
forms of the common hypothesis, not its unconditional proof.
At degree one, positivity says $0\le u_1\le1$; at degree two,
it additionally says $0\le u_2\le u_1$ and $u_2\ge2u_1-1$.
These are useful necessary conditions on the first moments, obtained
without analytic machinery. The full condition requires every such
finite-difference inequality at every order, so these low-degree
constraints do not replace its quantifiers.
The actual specializations are present in \texttt{RouteOne.lean},
including \texttt{H00\_iff\_all\_variation\_one} and
\texttt{H00\_iff\_actual\_weights\_nonnegative}~\cite{source}.

\PaperPage{19}{2.5 Finite averages and the shifted-moment recurrence}
\label{bern:recurrence-page}

For a function on the compact interval define the finite Bernstein
average using the explicit degree-zero convention
\[
 q_{N,j}=\frac jN\quad(N\ge1),\qquad q_{0,0}=0,
 \qquad A_N^u(f)=\sum_{j=0}^Nw^u_{N,j}f(q_{N,j}).
\]
Thus $A_0^u(f)=u_0f(0)$; a constant monomial always means the
constant function one, including its value at zero. No sign hypothesis
is imposed in this definition. Let $u^+_r=u_{r+1}$ be the shifted
sequence. Directly from~\eqref{bern:weights},
\begin{equation}\label{bern:shift}
 (j+1)w^u_{N+1,j+1}=(N+1)w^{u^+}_{N,j}
 \qquad(0\le j\le N).
\end{equation}
Indeed, the inner finite differences on the two sides are identical,
and their outer coefficients agree by
$(j+1)\binom{N+1}{j+1}=(N+1)\binom Nj$.

For a positive-degree monomial the $j=0$ contribution at degree
$N+1$ vanishes. Reindexing all remaining terms and using
\eqref{bern:shift} to absorb one power of the grid coordinate gives
\begin{equation}\label{bern:monomial-shift}
 A_{N+1}^u(x^{k+1})=
 \sum_{j=0}^Nw^{u^+}_{N,j}
       \left(\frac{j+1}{N+1}\right)^k.
\end{equation}
This formula also holds at $N=0$. To express its right-hand side
through averages at degree $N$, assume now $N\ge1$ and use
\[
 \frac{j+1}{N+1}=\frac{N}{N+1}\frac jN+\frac1{N+1},
 \qquad
 T_{N,k,i}=\binom ki\left(\frac N{N+1}\right)^i
                    \left(\frac1{N+1}\right)^{k-i}.
\]
The ordinary binomial theorem, followed by interchange of two finite
sums, proves the recurrence
\begin{equation}\label{bern:recurrence}
 A_{N+1}^u(x^{k+1})=
       \sum_{i=0}^kT_{N,k,i}A_N^{u^+}(x^i)
       \qquad(N\ge1).
\end{equation}
Its degree is lower on the right, but its moment sequence is shifted.
This matters for the induction on the next page: the inductive
statement must quantify over all sequences, not just a single fixed
one. The transition coefficients depend only on $N,k,i$. Although
they are nonnegative and sum to one, the averages multiplying them
can have either sign. Their convergence will come from induction,
without a bound on the variation of the original sequence.
There is also a useful exact first-moment check on the reindexing.
Putting $k=0$ in~\eqref{bern:monomial-shift} reduces its right-hand
side to the mass of a row for $u^+$, namely $u_1$. Thus every
positive-degree grid recovers the first moment exactly, even when
some of its weights are negative. The exceptional grid at degree
zero has average zero on $x$, and does not affect any limiting
statement. Stating that exception explicitly avoids silently applying
an affine-grid calculation with a zero denominator.

\PaperPage{20}{2.6 Monomial convergence without positivity}
\label{bern:convergence-page}

\begin{theorem}[Algebraic moment recovery]\label{bern:moment-limit}
For every real sequence $u$ and every integer $k\ge0$,
\begin{equation}\label{bern:limit}
 \lim_{N\to\infty}A_N^u(x^k)=u_k.
\end{equation}
\end{theorem}
\begin{proof}
We use strong induction on $k$, keeping the sequence universally
quantified. At degree zero, row mass gives the exact identity
\[
 A_N^u(1)=\sum_{j=0}^Nw^u_{N,j}=u_0
 \qquad\text{for every }N\ge0.
\]
Suppose the assertion holds in every degree at most $k$, for every
real sequence. In particular, it holds for the shifted sequence
$u^+$. The elementary scalar limits $N/(N+1)\to1$ and
$1/(N+1)\to0$ imply, for each $0\le i\le k$,
\[
 T_{N,k,i}\longrightarrow
 \begin{cases}1,&i=k,\\0,&i<k,\end{cases}
 \qquad A_N^{u^+}(x^i)\longrightarrow u_{i+1}.
\]
For $i<k$ the exponent $k-i$ is positive; for $i=k$ it is zero
and the binomial coefficient is one. Taking limits in the finite
sum~\eqref{bern:recurrence} therefore gives
\[
 A_{N+1}^u(x^{k+1})\longrightarrow
       \sum_{i=0}^k\mathbf1_{\{i=k\}}u_{i+1}=u_{k+1}.
\]
The recurrence is valid for every $N\ge1$, which is sufficient for
this limit. Removing the fixed index shift preserves convergence,
and proves the claim at degree $k+1$.
\end{proof}

The proof does not assume normalization, bounded moments, or
nonnegative weights. It does not exchange an infinite series with
a limit: for each fixed degree there are finitely many transition
terms. It also does not assert convergence for every continuous
test function when the weights have unrestricted signs. Such an
extension needs uniform control of the functionals, which positivity
will provide below.

Linearity immediately extends the conclusion to every polynomial:
\[
 p(x)=\sum_{k=0}^dc_kx^k
 \quad\Longrightarrow\quad
 A_N^u(p)\longrightarrow\sum_{k=0}^dc_ku_k=\mathcal M_u(p).
\]
As elementary checks, the recurrence gives
$A_N^u(x)=u_1$ for $N\ge1$ and
$A_N^u(x^2)=N^{-1}u_1+(1-N^{-1})u_2$ for $N\ge1$.
These are consistent both with degree one and with the limiting
second moment.

The absence of growth assumptions is possible because a fixed
polynomial degree probes only finitely many initial moments. The
recurrence reduces that degree while shifting the sequence, and the
induction terminates at row mass. It does not provide estimates
uniform in the polynomial degree. Consequently, using this theorem
inside an infinite power-series argument requires an additional
summability or domination argument; that step is kept separate in
the later resolvent analysis.

The general induction is the argument encoded by
\texttt{bernsteinAverage\_monomial\_tendsto} in
\texttt{BernsteinMomentLimit.lean}~\cite{source}.

\PaperPage{21}{2.7 A single ultrafilter and a positive functional}
\label{bern:functional-page}

Assume henceforth $u_0=1$ and $w^u_{N,j}\ge0$ for all admissible
indices. Each row is now a probability vector. On the real Banach
space $C([0,1])$, the finite averages obey
\[
 \begin{gathered}
 A_N^u(\alpha f+\beta g)=\alpha A_N^u(f)+\beta A_N^u(g),\\
 A_N^u(1)=1,\qquad f\ge0\Longrightarrow A_N^u(f)\ge0.
 \end{gathered}
\]
Moreover, positivity and mass one imply the common norm bound
\begin{equation}\label{bern:uniform-bound}
 |A_N^u(f)|\le\sum_{j=0}^Nw^u_{N,j}|f(q_{N,j})|
             \le\|f\|_\infty.
\end{equation}
Choose one free ultrafilter $\mathcal U$ on $\mathbb N$ containing
every cofinite set, and use this same choice for every test function.
For each $f$, the sequence of averages lies in the compact interval
$[-\|f\|_\infty,\|f\|_\infty]$. Compactness gives an ultrafilter
limit, and the Hausdorff property of the real line makes it unique.
Define
\[
 L_u(f)=\lim_{N\to\mathcal U}A_N^u(f).
\]
The use of a single ultrafilter is essential to this construction:
it lets addition and scalar multiplication pass to limits in a
common limiting process. An arbitrary, separately selected cluster
point for each function would not by itself establish linearity.

\begin{proposition}\label{bern:positive-functional}
The map $L_u:C([0,1])\to\mathbb R$ is a positive bounded linear
functional with $L_u(1)=1$ and operator norm one.
\end{proposition}
\begin{proof}
The finite linear identities pass to $\mathcal U$-limits because
addition and scalar multiplication are continuous. The closed sets
$[0,\infty)$ and $[-\|f\|_\infty,\|f\|_\infty]$ preserve,
respectively, positivity and the bound~\eqref{bern:uniform-bound}.
The constant sequence $A_N^u(1)=1$ gives $L_u(1)=1$. Therefore
$\|L_u\|\le1$, while evaluation on the unit vector $1$ gives the
reverse inequality.
\end{proof}

Because $\mathcal U$ contains every cofinite set, each ordinary
convergent sequence has the same ultrafilter limit. The preceding
page therefore yields
\[
 L_u(x^k)=u_k\qquad(k\ge0).
\]
The source uses \texttt{Filter.hyperfilter} for the common choice.

Its \texttt{hausdorffPositiveFunctional} is constructed from these limits. This is an
existence construction using classical choice; it supplies a coherent
functional before any representing measure is introduced.
It is useful to distinguish the two compactness inputs. Compactness of
$[0,1]$ makes each continuous test function bounded; the
resulting compact interval of real values supplies the ultralimit.
No subsequence indexed simultaneously for an uncountable set of
test functions is selected. The single ultrafilter provides that
simultaneous consistency through preservation of finite algebraic
operations.

\PaperPage{22}{2.8 The Hausdorff measure and uniqueness}
\label{bern:measure-page}

\begin{theorem}[Probability representation]\label{bern:measure}
If $u_0=1$ and every Bernstein weight of $u$ is nonnegative, there
is a unique Borel probability measure $\sigma$ on $[0,1]$ satisfying
\begin{equation}\label{bern:measure-moments}
 u_k=\int_{[0,1]}q^k\,d\sigma(q)\qquad(k\ge0).
\end{equation}
\end{theorem}
\begin{proof}
The real Riesz representation theorem on the compact Hausdorff
space $[0,1]$ applies to the positive functional $L_u$. It gives a
regular Borel measure with
\[
 L_u(f)=\int f\,d\sigma\quad(f\in C([0,1])),
 \qquad \sigma([0,1])=L_u(1)=1.
\]
Taking monomials and using $L_u(x^k)=u_k$ proves the moment
identities. For uniqueness, let $\sigma$ and $\tau$ be two Borel
probability measures with these moments. Their integrals agree on
polynomials by linearity. Polynomial density in $C([0,1])$ implies
that for each continuous $f$ and $\varepsilon>0$ there is a
polynomial $p$ with $\|f-p\|_\infty<\varepsilon$. Hence
\[
 \left|\int f\,d\sigma-\int f\,d\tau\right|
 \le\int|f-p|\,d\sigma+\int|f-p|\,d\tau<2\varepsilon.
\]
The continuous integrals therefore agree. Finite Borel measures on
this compact metric space are regular, and uniqueness in Riesz
representation gives $\sigma=\tau$.
\end{proof}

The converse is immediate but useful: any probability measure with
the displayed moments gives
\[
 w^u_{N,j}=\int B_{N,j}(q)\,d\sigma(q)\ge0.
\]
Thus normalized Bernstein positivity is exactly the compact-interval
moment condition. Uniqueness also removes any dependence of the
resulting measure on the free ultrafilter selected in the construction.
In fact the entire sequence of averages converges weakly to this
measure: polynomial approximation and~\eqref{bern:uniform-bound}
give $A_N^u(f)\to\int f\,d\sigma$ for every continuous $f$, by
bounding the two approximation errors and then taking the polynomial
limit.
This also proves directly that the limiting functional is intrinsic
to the moments: any bounded linear functional agreeing with it on
all monomials agrees on polynomials and then, by continuity, on all
continuous functions. The compact support is crucial to this density
argument. Moment uniqueness on a noncompact interval would require
additional hypotheses and is not being inferred here.

For the actual moments, Theorem~\ref{bern:h00} supplies positivity
from the same hypothesis $H_{00}$. The source construction
\texttt{actualHausdorffMeasure} and its moment theorem implement this
conditional measure step in \texttt{HausdorffMeasure.lean}~\cite{source}.
The uniqueness and weak convergence arguments just given are
mathematical exposition; their inclusion here does not claim a new,
separately audited Lean declaration. No unconditional proof of
$H_{00}$ has entered the construction.
